% F10：本书原创矢量示意；依据所在节的定义、证明或明确模型。
\begin{figure}[htbp]
\centering
\begin{tikzpicture}[x=1cm,y=1cm,>=stealth,every node/.style={font=\small},line width=.65pt]
\node[] at (6,3.5) {三分 Cantor 集的有限阶段与自然覆盖};
\draw[AnalysisBlue,line width=2pt] plot coordinates {(1.0000,2.6000) (11.0000,2.6000)};
\node[anchor=east] at (0.1,2.6) {\(j=0\)};
\node[anchor=west] at (12,2.6) {\(N=2^0\)};
\draw[AnalysisBlue,line width=2pt] plot coordinates {(1.0000,1.8500) (4.3333,1.8500)};
\draw[AnalysisBlue,line width=2pt] plot coordinates {(7.6667,1.8500) (11.0000,1.8500)};
\node[anchor=east] at (0.1,1.85) {\(j=1\)};
\node[anchor=west] at (12,1.85) {\(N=2^1\)};
\draw[AnalysisBlue,line width=2pt] plot coordinates {(1.0000,1.1000) (2.1111,1.1000)};
\draw[AnalysisBlue,line width=2pt] plot coordinates {(3.2222,1.1000) (4.3333,1.1000)};
\draw[AnalysisBlue,line width=2pt] plot coordinates {(7.6667,1.1000) (8.7778,1.1000)};
\draw[AnalysisBlue,line width=2pt] plot coordinates {(9.8889,1.1000) (11.0000,1.1000)};
\node[anchor=east] at (0.1,1.1) {\(j=2\)};
\node[anchor=west] at (12,1.1) {\(N=2^2\)};
\draw[AnalysisBlue,line width=2pt] plot coordinates {(1.0000,0.3500) (1.3704,0.3500)};
\draw[AnalysisBlue,line width=2pt] plot coordinates {(1.7407,0.3500) (2.1111,0.3500)};
\draw[AnalysisBlue,line width=2pt] plot coordinates {(3.2222,0.3500) (3.5926,0.3500)};
\draw[AnalysisBlue,line width=2pt] plot coordinates {(3.9630,0.3500) (4.3333,0.3500)};
\draw[AnalysisBlue,line width=2pt] plot coordinates {(7.6667,0.3500) (8.0370,0.3500)};
\draw[AnalysisBlue,line width=2pt] plot coordinates {(8.4074,0.3500) (8.7778,0.3500)};
\draw[AnalysisBlue,line width=2pt] plot coordinates {(9.8889,0.3500) (10.2593,0.3500)};
\draw[AnalysisBlue,line width=2pt] plot coordinates {(10.6296,0.3500) (11.0000,0.3500)};
\node[anchor=east] at (0.1,0.3500000000000001) {\(j=3\)};
\node[anchor=west] at (12,0.3500000000000001) {\(N=2^3\)};
\node[] at (6,-0.65) {\(r_j=3^{-j},\qquad N_jr_j^s=(2\cdot3^{-s})^j\)};

\end{tikzpicture}
\caption[覆盖尺度与 Hausdorff 维数]{覆盖尺度与 Hausdorff 维数。自然覆盖的代价在 \(s=\log2/\log3\) 两侧表现不同。图只展示这些覆盖的上界机制；维数的反向估计仍需质量分布等论证。}
\label{b1:fig:visual-f10}
\end{figure}
